\section{Experimental Methodology}
\label{sec:experiments}

\subsection{Models and Panel Construction}

We select 15 LLM judges spanning 6 provider families and 3 capability tiers, listed in Table~\ref{tab:models}. Nine models are accessed via the OpenRouter API; six open-weight models are served locally using vLLM~0.21.0 with AWQ INT4 quantization on an NVIDIA RTX~6000 Ada (48\,GB). API models are assigned to \emph{frontier}, \emph{mid-tier}, or \emph{lightweight} tiers following the provider capability classification; all six local models fall into the lightweight tier, as INT4 quantization substantially reduces judge accuracy on the pairwise judgment task (Table~\ref{tab:models}). We study three panel sizes $K \in \{3, 5, 7\}$ (all odd, majority-vote threshold $t = \lceil K/2 \rceil$). Exhaustive enumeration yields $\binom{15}{3} = 455$ panels at $K{=}3$, $\binom{15}{5} = 3{,}003$ at $K{=}5$, and $\binom{15}{7} = 6{,}435$ at $K{=}7$, covering every combination of families and tiers. Each panel evaluates all benchmark items under all attack types, producing a complete panel $\times$ item $\times$ attack tensor. The primary analysis uses $K{=}3$; \S\ref{sec:cross_k} extends to larger panels.

\begin{table}[t]
\centering
\small
\begin{tabular}{lllccc}
\toprule
\textbf{Model} & \textbf{Family} & \textbf{Tier} & \textbf{Deploy} & \textbf{MMLU} & \textbf{ARC} \\
\midrule
gpt-5.5              & OpenAI    & frontier    & API   & .957 & .983 \\
gemini-3.1-pro-preview & Google  & frontier    & API   & .955 & .987 \\
claude-opus-4-6       & Anthropic & frontier   & API   & .890 & .987 \\
gemini-3.5-flash      & Google    & frontier   & API   & .968 & 1.00$^\dagger$ \\
\addlinespace
gpt-4.1               & OpenAI    & mid-tier   & API   & .837 & .967 \\
claude-sonnet-4-6     & Anthropic & mid-tier   & API   & .810 & .953 \\
\addlinespace
gpt-4.1-mini          & OpenAI    & lightweight & API  & .800 & .963 \\
gpt-4.1-nano          & OpenAI    & lightweight & API  & .633 & .853 \\
claude-3-haiku        & Anthropic & lightweight & API  & .767 & .450 \\
mistral-large (123B)  & Mistral   & lightweight & local & .537 & .533 \\
llama3.1-70b          & Llama     & lightweight & local & .523 & .520 \\
llama3.1-8b           & Llama     & lightweight & local & .543 & .320 \\
qwen2.5-72b           & Qwen      & lightweight & local & .333 & .410 \\
qwen2.5-32b           & Qwen      & lightweight & local & .373 & .450 \\
qwen2.5-14b           & Qwen      & lightweight & local & .253 & .240 \\
\bottomrule
\end{tabular}
\caption{Judge models used in the study (clean accuracy on 300-item samples). API models accessed via OpenRouter (May 2026 snapshots); local models served with vLLM + AWQ INT4 quantization. $^\dagger$\,gemini-3.5-flash achieves perfect ARC accuracy but with 26.3\% format error rate. Tier assignment follows provider capability classification (frontier $>$85\% clean accuracy on both benchmarks, mid-tier 80--85\%, lightweight $<$80\% or quantized); all six local AWQ INT4 models fall into lightweight, where quantization-induced accuracy reduction conflates genuine capability limitations with quantization damage (\S\ref{sec:analysis}).}
\label{tab:models}
\end{table}

\subsection{Benchmarks and Attack Types}

Two benchmarks test distinct evaluation capabilities: MMLU \citep{hendrycks2021mmlu}, a knowledge-intensive multiple-choice benchmark spanning 57 subjects, and ARC \citep{clark2018arc}, a reasoning-intensive science question benchmark. We sample 300 items from each benchmark. MMLU probes factual recall and broad domain knowledge; ARC requires multi-step scientific reasoning. All API calls use temperature~0 for deterministic outputs; local models are served with the same temperature setting (a pilot study confirms temperature has negligible effect on $K_{\mathrm{eff}}$; Appendix~\ref{app:temperature}). Attacked variants are generated by programmatically applying each manipulation to the candidate text, with a fixed random seed for reproducibility.

Each benchmark is paired with six attack types that target different judge vulnerabilities: (1)~\textbf{prompt injection}, adversarial instructions embedded in candidate text that override the evaluation prompt \citep{perez2022ignore,greshake2023not}; (2)~\textbf{sycophancy}, responses engineered to agree with or flatter the judge \citep{wang2023fair}; (3)~\textbf{score manipulation}, explicit numerical cues (e.g., ``this answer deserves a 10/10'') injected into the candidate to anchor scoring; (4)~\textbf{verbosity bias}, padding correct answers with excessive irrelevant detail to exploit length preferences \citep{zheng2023judging}; (5)~\textbf{position bias}, reordering candidate answers to place the target in the position judges systematically favor \citep{wang2023fair}; and (6)~\textbf{authority bias}, prefacing candidate text with fabricated credentials or citations to inflate perceived credibility. The combination of 2 benchmarks and 6 attacks produces 12 experimental conditions, each evaluated across all 455 panels (implementation details and per-attack parameters in Appendix~\ref{app:attack_details}, Table~\ref{tab:attack-params}).

\subsection{Vulnerability Metrics}

Per-judge vulnerability under attack $a$ is the fraction of items where judge $j$ is fooled: $\eta_j(a) = \Pr(X_j^a {=} 1)$, estimated empirically over the item set. The weakest-link vulnerability $\eta_{\max}(a) = \max_{j \in \mathcal{P}} \eta_j(a)$ captures the most susceptible member of each panel. The panel attack success rate $\mathrm{ASR}(\mathcal{P}, a)$ is the fraction of items where at least $t$ judges are simultaneously fooled. Panel diversity is quantified by $K_{\mathrm{eff}}$ (Equation~\ref{eq:keff}), computed from pairwise mutual information between judges' three-category winner labels (model~A wins, model~B wins, tie) on clean (unattacked) data. MI is estimated via the plug-in estimator with Miller--Madow bias correction. The term $\exp(-2\,\mathrm{MI}(i,j))$ provides an information-theoretic measure of pairwise alienation: it equals 1 when judges are independent ($\mathrm{MI}{=}0$) and approaches 0 as dependence grows. Under a Gaussian copula, this reduces to $1 - \rho_{ij}^2$, yielding the closed form $K_{\mathrm{eff}} = K - \frac{1}{K}\sum_{i}\sum_{j \neq i} \rho_{ij}^2$ \citep{nelsen2006copulas}; however, $K_{\mathrm{eff}}$'s validity does not depend on this parametric equivalence. Goodness-of-fit testing rejects the Gaussian copula for 97--100\% of judge pairs (Appendix~\ref{app:copula_gof}), yet the plug-in MI and closed-form rankings agree closely (within-$K$ Spearman $\rho{\approx}0.86$), confirming that $K_{\mathrm{eff}}$ is robust to copula misspecification as a ranking metric. Together, $\eta_{\max}$ and $K_{\mathrm{eff}}$ serve as the two predictors in the regression framework below: $\eta_{\max}$ indexes the weakest-link channel and $K_{\mathrm{eff}}$ indexes the diversity channel. Robustness of $\eta_{\max}$ under partial panel information is analyzed in Appendix~\ref{app:oracle}.

\subsection{Aggregation Methods}

Beyond the default majority vote, we evaluate three alternative aggregation rules to test whether the diversity-vulnerability association persists under different decision mechanisms. The aggregation comparison uses the same log-log regression specification as the primary analysis (Equation~\ref{eq:regression}).

\paragraph{Majority Vote.} Accept if $\sum_{j=1}^K \mathbf{1}[J_j \text{ approves}] \geq \lceil K/2 \rceil$ (baseline).

\paragraph{Accuracy-Weighted Vote.} Weight judges by clean-data accuracy: $w_j \propto \mathrm{acc}_j$, $\sum_j w_j = 1$; accept if $\sum_j w_j \cdot \mathbf{1}[J_j \text{ approves}] > 0.5$.

\paragraph{Threshold Vote (Unanimity).} Accept only if all $K$ judges approve ($t = K$).

\paragraph{Bayesian Calibration.} Compute $P(\text{correct} \mid \mathbf{v}) \propto P(\text{correct}) \prod_j P(v_j \mid \text{correct})$ using each judge's confusion matrix on clean data; accept if the posterior exceeds 0.5.

\subsection{Statistical Framework}
\label{sec:stat_framework}

To disentangle the contributions of panel diversity and individual judge robustness, we fit a log-log OLS regression for each of the 12 conditions:
\begin{equation}
\label{eq:regression}
\log(\mathrm{ASR} + \epsilon) = \alpha + \beta_1 \log(K_{\mathrm{eff}}) + \beta_2 \log(\eta_{\max} + \epsilon),
\end{equation}
with $\epsilon = 10^{-6}$ to handle zero values (sensitivity to $\epsilon$ is analyzed in Appendix~\ref{app:epsilon}). Standardized coefficients $\hat{\beta}_{\mathrm{std}} = \hat{\beta} \cdot \mathrm{SD}(\log x) / \mathrm{SD}(\log y)$ enable direct comparison of effect magnitudes across predictors and conditions. We apply Benjamini-Hochberg correction for multiple testing, controlling the false discovery rate per predictor across the 12 conditions.

To assess relative predictor importance beyond regression coefficients, we perform dominance analysis \citep{budescu1993dominance}. This method decomposes $R^2$ by averaging each predictor's incremental $R^2$ contribution across all possible subsets of predictors. A predictor is \emph{dominant} if its average contribution exceeds that of the competing predictor in every subset ordering. With two predictors, dominance reduces to comparing each predictor's solo $R^2$ with its incremental $R^2$ conditional on the other.

To test whether $K_{\mathrm{eff}}$'s effect direction varies by attack type, we fit generalized additive models (GAMs) with smooth terms for both predictors, allowing nonlinear partial effects. The GAM partial effect of $K_{\mathrm{eff}}$ is then examined separately for targeted attacks (prompt injection, score manipulation, position bias) and shared-bias attacks (sycophancy, verbosity bias, authority bias).

\subsection{Robustness to Non-Independence}
\label{sec:robustness_method}

The 455 panels share judges combinatorially (each appears in 91 panels), inducing correlated residuals. We correct for this using clustered standard errors (CR1, 3-position max-SE), mixed-effects models (random judge intercepts, 3-position max-$p$), and wild cluster bootstrap ($B{=}9{,}999$, Rademacher weights, 3-position max-$p$) following \citet{cameron2008bootstrap}. The 13 clusters per position at $K{=}3$ are below the ${\geq}20$ threshold recommended by \citet{cameron2008bootstrap}; at $K{=}5$, cluster count drops to ${\leq}11$, and at $K{=}7$ to ${\leq}9$. Following \citet{mackinnon2018wild}, we use Rademacher weights, which remain conservative in few-cluster settings even when cluster sizes are unequal. The bootstrap attenuation at $K{=}3$ (4/12 significant) may partly reflect insufficient cluster count rather than absent effects. We use the log-log specification rather than fractional response models \citep{papke1996econometric} or beta regression \citep{ferrari2004beta} because the log transform linearizes the multiplicative $K_{\mathrm{eff}} \cdot \eta_{\max}$ structure; fractional logit with clustered SE serves as an $\epsilon$-free robustness check (Appendix~\ref{app:robustness_checks}). Details are in Appendix~\ref{app:robustness_methods}.
